\documentclass[11pt]{article}
\usepackage[margin=1.2in]{geometry}
\usepackage{amsmath,amssymb,amsthm}
\usepackage{hyperref}
\hypersetup{colorlinks=true, linkcolor=blue, urlcolor=blue}
\usepackage{parskip}
\emergencystretch=3em

\newcommand{\R}{\mathbb{R}}
\newcommand{\code}[1]{\texttt{#1}}

\title{Tide retrospective: the radial Taylor bound\\
\large grand composition, final}
\author{Claude (Anthropic), with GPT-5.6 Sol consults}
\date{2026-08-10}

\begin{document}
\maketitle

\begin{abstract}
The last open Lean item of the inverse-direction programme. A
globally $C^k$ loss satisfies, on every ball, the order-$k$
fixed-ball Taylor remainder bound that the package structures carry
as a field (\code{exists\_taylorRemainder\_bound}), by the radial
route: the ray-derivative identity at arbitrary parameter values
(generalizing the corpus's origin-only lemma, whose composition proof
was already global in the point), the one-dimensional
Taylor--Lagrange theorem along each ray, the compactness bound for
the $k$-th Fr\'echet derivative on the closed ball, and the
multilinear operator-norm estimate, giving $C = M/k!$. The corollary
\code{higherLaplaceDomainFamily\_ofContDiff} discharges the
package-derivation wrapper's remaining hypothesis, and with it the
inverse Theorem 3.1 chain is hypothesis-complete from the note's
prose data: smoothness, a vanishing gradient, and the
diagonal-matched positive-definite matrix produce the certified
family; analyticity at the minima and one compactly-supported
physical data premise produce the location and the germ. The scoping
consult rated this tide high-risk and API-sensitive; it compiled
with zero fix rounds, because each flagged friction point had an
exact Mathlib counterpart.
\end{abstract}

\section{Setting}

The package-derivation wrapper assumed its Taylor bounds; the
consult's Rank-2 plan --- a distinct engineering task, 70--160
lines, ``High/API-sensitive'' --- was to derive them from
smoothness. This tide executes that plan unchanged.

\section{The thing formalised}

\begin{sloppypar}
\code{ray\_iteratedDeriv\_at} adapts the corpus proof: the ray is
the composition of $L$ with \code{toSpanSingleton}, the iterated
Fr\'echet derivative of a composition with a continuous linear map
evaluates by \code{compContinuousLinearMap}, and the singleton map
sends $1$ to $x$ --- nothing in that chain cares that the base point
is the origin. The main theorem fixes $y$ in the ball, restricts $L$
to the ray $g(s) = L(s \cdot y)$, and runs Mathlib's
\code{taylor\_mean\_remainder\_lagrange} on $[0,1]$ at order
$k - 1$. The two bridging concerns the consult flagged dissolved on
contact: \code{iteratedDerivWithin\_eq\_iteratedDeriv} takes a
\code{ContDiffAt} hypothesis and holds at the interval's endpoint,
and \code{taylor\_within\_apply} is already the range-sum with
inverse-factorial coefficients, so the identification with
\code{taylorHomogeneousTerm} is the ray identity at $0$ plus
\code{ring}. The mean-point value is bounded by the closed-ball
supremum of the $k$-th derivative norm --- compactness plus
\code{ContDiff.continuous\_iteratedFDeriv} --- through the
multilinear operator norm with a constant product.
\end{sloppypar}

\section{Where progress stalled}

It did not: zero fix rounds on a tide the deliberation rated
high-risk. The credit belongs half to the consult (whose friction
list dictated exactly which bridging lemmas to look up before
writing) and half to the accumulated gotcha catalogue (the final
division monotonicity went to \code{gcongr} on the first attempt
rather than a named-lemma hunt).

\section{Mathlib lookup versus new mathematics}

Lookup:
\begin{itemize}
\item \code{taylor\_mean\_remainder\_lagrange} and
  \code{taylor\_within\_apply};
\item \code{iteratedDerivWithin\_eq\_iteratedDeriv};
\item the continuity and differentiability of iterated derivatives
  from \code{ContDiff};
\item \code{IsCompact.exists\_bound\_of\_continuousOn};
\item the multilinear operator-norm bound (\code{le\_opNorm}) and
  the composition rule
  \code{iteratedFDeriv\_comp\_right}.
\end{itemize}
New mathematics: none. The tide is the textbook multivariate Taylor
estimate, assembled from one-dimensional parts.

\section{What the next tide inherits}

Nothing mandatory --- the inverse-direction programme's Lean side is
closed. Optional niceties recorded in the tide log: a single
end-to-end convenience statement composing the family constructor
with the located analytic capstone, the note-side markers for the
new theorems (a gated TeX action), and the two review batches of
recorded proof refactors for a simplification pass.
\end{document}
